\chapter{Random Walks}
\label{ch:2}
\section{Paths and path integrals}
\label{sec:2-1}
Consider length-$N$ binary paths $w=(w_1,\ldots,w_N)$, where
$w_a\in\{0,1\}$ and $|w|=\sum_aw_a$ is the endpoint coordinate.
Fix $j$ distinguished steps and assign the signed weight
\[
\omega_j(w)=(-1)^{w_1+\cdots+w_j}.
\]
The particular first-$j$ choice is inessential; any fixed subset of $j$
steps gives the same endpoint sums. A path with endpoint $n$ and $r$
right steps among the distinguished steps has weight $(-1)^r$.
There are $\binom jr\binom{N-j}{n-r}$ such paths, so
\begin{equation}\label{eq:signed-path-Krawtchouk}
\sum_{|w|=n}\omega_j(w)
=\sum_r(-1)^r\binom jr\binom{N-j}{n-r}
=K_n(j;N).
\end{equation}
This is a path interpretation of the coefficient formula in
Section~\ref{sec:1-2}, not a new definition. For $j=0$ it counts ordinary
paths and gives $\binom Nn$; for $j>0$, cancellation produces signed
entries of column $j$ of $\Phi_N$. For instance, $N=4$, $j=2$ gives
endpoint sums $(1,0,-2,0,1)$.

The signed Galton-board construction of \cite[Section~6]{KocikFeynman}
is an elementary analogue of a sum over paths. Its weights are not
probabilities. In quantum mechanics one sums amplitudes before taking
an appropriately normalized absolute square; a physical unitary model
requires more structure than this combinatorial identity. The unsigned
walk and the stochastic Ehrenfest transition operator should likewise
not be identified with a unitary evolution.


\subsection{Weighted paths and the two-parameter deformation}
\label{sec:weighted-path-deformation}
Over a commutative ring, define a separate coefficient array
\[
\mathcal K_n(j;N;\alpha,\beta)
=[v^n](1+\alpha v)^{N-j}(1+\beta v)^j.
\]
A consistent path rule gives every left step weight $1$, and a right
step weight $\beta$ at each of $j$ distinguished times and $\alpha$ at
all other times. Thus
\[
\mathcal K_n(j;N;\alpha,\beta)
=\sum_r\binom jr\binom{N-j}{n-r}\beta^r\alpha^{n-r}.
\]
This recovers $K_n(j;N)$ at $(\alpha,\beta)=(1,-1)$ without changing
its normalization. Coefficient comparison gives
\begin{align*}
\mathcal K_n(j;N+1)&=\mathcal K_n(j;N)+\alpha\mathcal K_{n-1}(j;N),\\
\mathcal K_n(j+1;N+1)&=\mathcal K_n(j;N)+\beta\mathcal K_{n-1}(j;N).
\end{align*}
For adjacent entries $\begin{pmatrix}a&b\\c&d\end{pmatrix}$ in this
array, $\beta a+c=\alpha b+d$.
These identities require no division by $\alpha-\beta$ and remain valid
when parameters coincide or are not invertible.

\begin{remark}[A discrepancy in the source's path rule]
The displayed generating function of \cite[Section~6.2]{KocikFeynman}
is the one used here. However, the preceding prose assigns left/right
weights $\alpha,\beta$ only at the $j$ distinguished times, and weights
$1,1$ elsewhere. That different rule produces
$(1+v)^{N-j}(\alpha+\beta v)^j$, not the displayed generating function.
We retain the displayed coefficient array and specify the compatible
rule explicitly. No stochastic normalization, orthogonality, or
involution property is asserted for arbitrary parameters.
\end{remark}

\section{$A$, $K$, and $\Lambda$}
\label{sec:2-2}
For the symmetric walk on $\{0,\ldots,N\}$, use column vectors for
state distributions and the zero-based standard basis $e_j$. Put
\[
A_Ne_j=(N-j)e_{j+1}+je_{j-1},\qquad
\Lambda_N=\operatorname{diag}(N,N-2,\ldots,-N),
\]
with endpoint terms omitted. In this section the matrix symbol $K$
is explicitly defined to mean $\Phi_N=(K_n(j;N))_{n,j=0}^N$.
It is not a single polynomial $K_n$ or the symmetric matrix
$S_N=\Phi_NB_N$.

\subsection{Intertwining and diagonalization}
Writing $G_j(v)=(1+v)^{N-j}(1-v)^j$, differentiation gives
\[
(1-v^2)G_j'(v)=(N-2j-Nv)G_j(v).
\]
Equating coefficients yields
\[
(N-n+1)K_{n-1}(j;N)+(n+1)K_{n+1}(j;N)
=(N-2j)K_n(j;N).
\]
The spatial recurrence is
\[
jK_n(j-1;N)+(N-j)K_n(j+1;N)=(N-2n)K_n(j;N).
\]
Absent endpoint terms are zero. In matrix form these identities read
\[
A_N\Phi_N=\Phi_N\Lambda_N,\qquad
\Phi_NA_N=\Lambda_N\Phi_N.
\]
This is the coefficient proof of the master equation; its symmetric-power
interpretation is given in Section~\ref{sec:8-1}.
Since $\Phi_N^2=2^NI$, we obtain
\[
\Phi_N^{-1}=2^{-N}\Phi_N,\qquad
A_N=2^{-N}\Phi_N\Lambda_N\Phi_N.
\]
Column $j$ of $\Phi_N$ is a right eigenvector with eigenvalue $N-2j$;
row $n$ is a left eigenvector with eigenvalue $N-2n$.

\subsection{Transition probabilities}
For $N\geq1$, $T_N=A_N/N$ is column-stochastic. Hence
$p_{t+1}=T_Np_t$ and, for integers $t\geq0$,
\[
T_N^t=2^{-N}\Phi_N(\Lambda_N/N)^t\Phi_N.
\]
The backward operator on column vectors of function values is $T_N^T$,
not $T_N$; its eigenfunctions are $j\mapsto K_n(j;N)$.
The stationary column distribution is
$2^{-N}(\binom N0,\ldots,\binom NN)^T$, the normalized first column
of $\Phi_N$. The eigenvalue $-1$ of $T_N$ expresses period two;
stationarity does not imply convergence from arbitrary initial data.

\subsection{The cases $N=3$ and $N=5$}
For $N=3$,
\begingroup
\small
\setlength{\arraycolsep}{3pt}
\[
A_{3}=\begin{pmatrix}
0 & 1 & 0 & 0 \\
3 & 0 & 2 & 0 \\
0 & 2 & 0 & 3 \\
0 & 0 & 1 & 0
\end{pmatrix}.
\]
\endgroup
\begingroup
\small
\setlength{\arraycolsep}{3pt}
\[
\Lambda_{3}=\begin{pmatrix}
3 & 0 & 0 & 0 \\
0 & 1 & 0 & 0 \\
0 & 0 & -1 & 0 \\
0 & 0 & 0 & -3
\end{pmatrix}.
\]
\endgroup
The transition eigenvalues are $1,\frac{1}{3},- \frac{1}{3},-1$.
For $N=5$,
\begingroup
\small
\setlength{\arraycolsep}{3pt}
\[
A_{5}=\begin{pmatrix}
0 & 1 & 0 & 0 & 0 & 0 \\
5 & 0 & 2 & 0 & 0 & 0 \\
0 & 4 & 0 & 3 & 0 & 0 \\
0 & 0 & 3 & 0 & 4 & 0 \\
0 & 0 & 0 & 2 & 0 & 5 \\
0 & 0 & 0 & 0 & 1 & 0
\end{pmatrix}.
\]
\endgroup
\begingroup
\small
\setlength{\arraycolsep}{3pt}
\[
\Lambda_{5}=\begin{pmatrix}
5 & 0 & 0 & 0 & 0 & 0 \\
0 & 3 & 0 & 0 & 0 & 0 \\
0 & 0 & 1 & 0 & 0 & 0 \\
0 & 0 & 0 & -1 & 0 & 0 \\
0 & 0 & 0 & 0 & -3 & 0 \\
0 & 0 & 0 & 0 & 0 & -5
\end{pmatrix}.
\]
\endgroup
The transition eigenvalues are $1,\frac{3}{5},\frac{1}{5},- \frac{1}{5},- \frac{3}{5},-1$.
The corresponding $R_N$, $L_N$, $\Phi_N$ and $S_N$ tables are in
Appendix~\ref{app:matrix-tables}. Centered and lattice polynomial tables
are in Appendix~\ref{app:polynomial-tables}.

\section{Nonsymmetric walks}
\label{sec:2-3}
% TODO: Develop this section; add verified references, examples, and exercises.

\section{Symmetric Krawtchouk matrices and binomial expectations}
\label{sec:2-4}
% TODO: Develop this section; add verified references, examples, and exercises.

